% !PRINTED 147
\weberpage{147}
\begin{center}
{\large \S.\ 45.}

\vspace{0.5em}
{\large\gesperrt{Zweiter Beweis des Satzes von den symmetrischen Functionen.}}
\end{center}

Es sei $S$ eine ganze symmetrische Function der Variablen $\alpha_1,\alpha_2\ldots\alpha_n$. Die einzelnen Glieder dieser Function sind alle von der Form
\[
 M\alpha_1^{\nu_1}\alpha_2^{\nu_2}\ldots\alpha_n^{\nu_n},
\]
worin $M$ ein von den $\alpha$ unabhängiger Coëfficient ist. Diese Glieder sollen nun in bestimmter Weise angeordnet werden. Es soll, nachdem die Reihenfolge der Variablen $\alpha_1,\alpha_2\ldots\alpha_n$ festgesetzt ist, von zwei Gliedern
\[
 A=M\alpha_1^{\nu_1}\alpha_2^{\nu_2}\ldots\alpha_n^{\nu_n},\qquad
 A'=M'\alpha_1^{\nu_1}\alpha_2^{\nu'_2}\ldots\alpha_n^{\nu'_n}
\]
% [sic] Source p147 omits the prime in the first exponent of A'; retained, not silently normalized.
$A$ das \gesperrt{höhere} genannt werden, wenn die erste der Differenzen\begingroup\renewcommand{\thefootnote}{E1}\footnote[1]{[Quellenhinweis:] Im Druck steht im ersten Exponenten des zweiten Gliedes $\nu_1$ ohne Strich. Die nachfolgende Definition erfordert dort $\nu'_1$; der gedruckte Ausdruck bleibt unverändert.}\endgroup
\[
 \nu_1-\nu'_1,\nu_2-\nu'_2\ldots\nu_n-\nu'_n,
\]
die von Null verschieden ist, einen \gesperrt{positiven} Werth hat, wenn also entweder $\nu_1>\nu'_1$ oder $\nu_1=\nu'_1,\nu_2>\nu'_2$ oder $\nu_1=\nu'_1,\nu_2=\nu'_2,\nu_3>\nu'_3$ etc.

Da wir alle Glieder, in denen sämmtliche Exponenten $\nu_1,\nu_2\ldots\nu_n$ übereinstimmen, in ein Glied vereinigt voraussetzen, so ist hiernach von je zwei Gliedern entschieden, welches das höhere ist, und wenn $A$ höher als $A'$, $A'$ höher als $A''$ ist, so ist auch $A$ höher als $A''$.

Ist nun nach dieser Anordnung
\[
 A=M\alpha_1^{\nu_1}\alpha_2^{\nu_2}\ldots\alpha_n^{\nu_n}
\]
das höchste Glied unserer Function $S$, so folgt, dass
\[
 \nu_1\eqslantgtr\nu_2
\]
sein muss; denn wäre $\nu_1<\nu_2$, so würde das Glied
\[
 M\alpha_1^{\nu_2}\alpha_2^{\nu_1}\ldots\alpha_n^{\nu_n},
\]
das wegen der Symmetrie auch in $S$ vorkommen muss, in der Ordnung höher stehen als $A$, gegen die Voraussetzung. Ebenso folgt, dass $\nu_2\eqslantgtr\nu_3$ sein muss, denn wäre $\nu_2<\nu_3$, so würde
\[
 M\alpha_1^{\nu_1}\alpha_2^{\nu_3}\alpha_3^{\nu_2}\ldots\alpha_n^{\nu_n}
\]
höher stehen als $A$ u.\ s.\ f. Wir schliessen also, dass die Exponenten in $A$
%CONT>>
